\documentclass[10pt,a4paper]{article}
\usepackage[margin=1.5cm]{geometry}
\usepackage{amsmath,amssymb}
\usepackage{enumitem}
\usepackage{titlesec}
\usepackage{hyperref}
\usepackage{xcolor}

\pagestyle{plain}
\setlength{\parindent}{0pt}
\setlength{\parskip}{4pt}

\titleformat{\section}{\bfseries\Large}{}{0em}{}[\titlerule]
\titleformat{\subsection}{\bfseries\large}{}{0em}{}
\titlespacing{\section}{0pt}{10pt}{6pt}
\titlespacing{\subsection}{0pt}{6pt}{3pt}

\title{\textbf{COSC 419 --- Study Notes}}
\author{}
\date{}

\begin{document}
\maketitle
\tableofcontents
\newpage

% ==============================================================
\section{Foundational Regression and Classification}
% ==============================================================

\textbf{Linear Regression vs.\ Logistic Regression:}
Linear regression is linear. Logistic regression is a sigmoid classifier.

% ==============================================================
\section{Loss Functions}
% ==============================================================

\subsection{Step Loss}
Bad to use because it is not differentiable.

\subsection{MSE (Mean Squared Error)}
Doesn't penalize wrong classification answers well enough (does penalize outlier numerical answers decently).

\subsection{Sigmoid}
An activation function that squashes inputs into a range between 0 and 1, often used in binary classification. It is prone to saturation, which causes vanishing gradients in deep networks.

\subsection{Cross-Entropy}
A log-based loss function that penalizes incorrect classifications harshly, with the cost approaching infinity as the predicted probability for the true class approaches zero.

\subsection{Relation between Cross-Entropy and Softmax}
Softmax is the activation function that turns raw scores (logits) into a probability distribution summing to 1. Cross-Entropy is the loss function that measures the distance between that distribution and the one-hot encoded true label.

% ==============================================================
\section{Regularization Techniques}
% ==============================================================

\subsection{L1 Regularization}
Adds the sum of absolute weights to the loss; it promotes sparsity, effectively performing feature selection by shrinking unimportant weights to exactly zero.

\subsection{L2 Regularization (Weight Decay)}
Adds the sum of squared weights to the loss; it prefers spreading weight evenly across all features, which improves stability and handles correlated features better than L1.

\subsection{Elastic Net}
Combines L1 and L2 using separate constants to balance sparsity and stability.

\subsection{Dropout}
Randomly sets neurons to zero each forward pass. Forces redundant learning.

\subsection{Batch Normalization}
Normalizes layer inputs using mini-batch mean and variance. It stabilizes training, allows for higher learning rates, and reduces internal covariate shift.

% ==============================================================
\section{Optimization Algorithms}
% ==============================================================

\subsection{Full-batch Gradient Descent}
Uses the entire dataset to compute one gradient update. Smooth loss and expensive.

\subsection{Mini-batch SGD}
Computes updates using a subset of the data, and updates based on a noisy estimate of the true gradient.

\subsection{SGD + Momentum}
SGD that has ``momentum'' that lets it move over some local minima.

\subsection{AdaGrad}
An adaptive algorithm that scales the learning rate for each parameter based on historical squared gradients.

\subsection{RMSProp}
A variation of AdaGrad that uses a moving average of squared gradients to prevent the learning rate from shrinking too quickly.

\subsection{Adam}
Combination of Momentum and RMSProp; it is often faster than standard SGD and handles sparse gradients well.

\subsection{AdamW}
Decouples weight decay from gradient updates in Adam.

% ==============================================================
\section{Multi-Layer Perceptrons (MLP)}
% ==============================================================

\textbf{Feature Transforms:} The process of projecting data into a new geometric space where previously inseparable data becomes linearly separable.

\textbf{Non-linearity:} Activation functions like ReLU or Sigmoid stop it from being linear, letting it learn more.

\textbf{Forward Pass:} Applying the weights in order through the neural network to make a prediction.

\textbf{Backpropagation:} Applying the chain rule to compute the partial derivative of the loss with respect to every weight, allowing for model updates.

All nodes are connected on each layer. Lots of sources means overfitting is a big risk.

% ==============================================================
\section{Convolutional Neural Networks (CNN)}
% ==============================================================

\subsection{CNN vs.\ MLP}
CNNs use weight sharing and locality (receptive fields) to handle spatial data efficiently, whereas MLPs are fully connected and require a unique parameter for every single input-output connection.

\subsection{Different Components}
Includes Convolutional layers for feature extraction, Pooling layers (Max/Average) for downsampling, and Fully Connected layers for the final classification.

\subsection{Calculations of Tensor Dimensions}
The size of a feature map is calculated as:
\[
\text{Output} = \left\lfloor \frac{W - F + 2P}{S} \right\rfloor + 1
\]
where $W$ is input size, $F$ is filter size, $P$ is padding, and $S$ is stride.

% ==============================================================
\section{CNN Architectures}
% ==============================================================

\begin{itemize}[nosep]
  \item \textbf{AlexNet} (2012) --- ReLU over tanh, 6$\times$ faster training
  \item \textbf{VGG} (2014) --- 3$\times$3 filters, more depth (19 layers)
  \item \textbf{GoogLeNet} (2014) --- Inception modules with parallel paths
  \item \textbf{ResNet} (2015) --- Introduced Skip Connections to solve vanishing gradients in very deep (152-layer) networks
  \item \textbf{ConvNeXt} (2022) --- Outcompetes transformers sometimes
\end{itemize}

% ==============================================================
\section{Training Your Neural Network}
% ==============================================================

\subsection{Activation Functions}
ReLU is the default for CNNs; Leaky ReLU fixes ``Dying ReLU'' (neurons stuck at 0); Sigmoid is reserved for output layers in binary tasks. ReLU is good because it is ``jagged'' across time.

\subsection{Weight Initialization}
He initialization is used for ReLU networks, double the variance of Xavier to make up for ReLU jaggedness; Xavier initialization is used for Sigmoid/Tanh to keep activation scales stable.

\subsection{Choosing Hyperparameters}
Key ``knobs'' include learning rate (most important), batch size, and weight decay (L2 regularization).

% ==============================================================
\section{Recurrent Neural Networks (RNN)}
% ==============================================================

\subsection{Language Modeling}
The task of predicting the next word or character in a sequence based on previous context.

\subsection{Sequence-to-Sequence}
An encoder-decoder framework used for tasks like translation, where one RNN compresses the input and another generates the output.

\subsection{RNN vs.\ LSTM}
Standard RNNs struggle with vanishing gradients over long sequences. LSTMs solve this using a Cell State (long-term memory highway) and gates (Forget, Input, Output) to regulate information flow.

\end{document}
